\section{Method}
\label{sec:method}

\subsection{Data}
\label{sec:method-data}

The two-moons dataset is used here; it is a binary classification problem where each
point consists of two features, and the classes are formed by
interlocking semicircles. The dataset requires a nonlinear, curved decision
boundary, can be completely visualized, and keeps the networks small enough to
make them solvable via MILP. A total of 2,400
points are created with Gaussian noise of 0.15 standard deviation (seed 1001) and
split into training (1,440), validation (480), and test (480) sets with equal proportion from each
class (seed 1002). There is a fixed scale that is set to be equal on both axes and is used to map
the square of half-width 2.25 centered at $(0.5, 0.25)$ to the input domain
$\D = [0,1]^2$; generation fails if any point is out of bounds. A list of file
hashes is used to verify that all computing platforms use the same data.

\subsection{Networks and training}
\label{sec:method-networks}

All networks considered are fully connected ReLU networks that produce two output scores. The
three architectures differ in depth and width: (a) a single hidden layer with 36
neurons (182 parameters, 36 ReLUs); (b) two hidden layers with 11 neurons each (189 parameters, 22
ReLUs); and (c) three hidden layers with 8 neurons each (186 parameters, 24 ReLUs). The parameter
count is balanced within 4 percent across architectures; however,
the number of ReLUs depends on the depth. As the certificate relaxes all uncertain
ReLUs, this difference represents a limitation.

Each architecture is trained with five random seeds (from 1 to 5), resulting in 15
trained networks. For the pilot experiment, seeds 9001 and 9002 are used. The training procedure
minimizes cross-entropy loss using the Adam optimizer, with a learning rate of
0.001 and a batch size of 64, for 500 epochs on one CPU thread in 32-bit floating point arithmetic.
All further evaluations are done in 64-bit arithmetic. The saved checkpoint is the one with the
lowest validation loss (in case of a tie, the earliest epoch is chosen); the test set is not
involved in the checkpoint choice process.

\subsection{Robustness radius}
\label{sec:method-radius}

The robustness radius of an input $x$ is
\begin{equation}
  \rstar(x) = \inf \left\{ \linf{x' - x} : x' \in \D,\ h(x') \neq h(x) \right\}
  \label{eq:radius}
\end{equation}
where $\linf{\cdot}$ is the L-infinity distance, i.e., the maximal amount of change that can be
made to any single coordinate, with the perturbed input $x'$ staying inside the domain $\D$. The
output $h(x)$ of a model is defined as the class with the higher score, where ties
are broken in favor of the class with the smaller index. For the inputs from the lower-indexed
class, a change of output requires the other class to score strictly higher; this
cannot be expressed using MILP, and therefore the substitute requirement for
Equation~\eqref{eq:radius} is an open problem: either the score boundary with a
given tolerance or a small required margin. In a first trial, the solver generated
a point with a score gap of $-5 \times 10^{-16}$, on the boundary
up to rounding but not a changed prediction. The rule will be fixed before the protocol is
frozen.

\subsection{Three estimates}
\label{sec:method-estimates}

The certificate provides a lower bound $L$. It uses the triangle LP
relaxation~\cite{MapSource09}, which replaces each uncertain ReLU by the minimal
convex region that is guaranteed to contain the graph of the ReLU inside the specified bounds, with
the interlayer bounds determined by interval arithmetic. Linear programs are solved by HiGHS
through SciPy, while a bisection method finds the maximal radius the relaxation can certify with
a precision of 0.0001. Since the relaxation contains all behavior of the network,
anything certified by the relaxation is guaranteed to be safe, i.e., $L \le \rstar$, but
the relaxation also allows impossible outputs, so $L$ can be
smaller than $\rstar$.

The attack gives an upper bound $U$. It is an untargeted L-infinity FAB~\cite{MapSource05},
performed with 5 restarts of 100 iterations with seed 3001; the concrete implementation has not yet
been chosen. A candidate is accepted only if it is contained in $\D$, changes the
prediction, and has a distance that matches an
independent recomputation to within $10^{-9}$. If such a candidate is found, then $\rstar \le U$
holds. The failure of the attack does not give any $U$, as not finding an
adversarial example does not mean there is none~\cite{MapSource03}.

The reference is the exact MILP~\cite{MapSource12}. The MILP has an on/off variable associated
with each uncertain ReLU. It is solved by HiGHS through SciPy on a single thread
with a 60-second limit, without initialization from the attack solution. The
solver outputs an interval $a \le \rstar \le b$, which is considered resolved
when $b - a \le 0.0001$.

\subsection{Error decomposition}
\label{sec:method-decomposition}

\begin{equation}
  \underbrace{U - L}_{\text{total gap}}
  = \underbrace{(\rstar - L)}_{\text{certificate error}}
  + \underbrace{(U - \rstar)}_{\text{attack error}}
  \label{eq:decomposition}
\end{equation}

Equation~\eqref{eq:decomposition} splits the total gap into two non-negative parts.
The certificate error $\rstar - L$ is the degree to which the certificate underestimates the
radius, while the attack error $U - \rstar$ is the degree to which the attack overestimates it.
Since the reference is an interval, the errors are bounded, not
exact. With the combined bounds $A = \max(a, L)$ and
$B = \min(b, U)$, the certificate error is within $[A - L,\ B - L]$ and the attack error is within
$[U - B,\ U - A]$. If $A > B$ by more than $10^{-6}$, then the estimates are inconsistent: the
instance is marked as inconsistent with no error recorded, and in the pilot, it is
considered a bug. A failed attack is counted without an attack error. Errors are
kept for each individual input and aggregated using medians, because a simple average may hide a
large miss on a single input~\cite{Carlini2017Provable}.

\subsection{Data collection and analysis}
\label{sec:method-collection}

The same 30 test inputs are used to evaluate all networks. These include 15
inputs from each class. They have been selected randomly with a fixed seed (seed
2001) before the computation of any robustness results. In total, 450 network--input instances (15
networks $\times$ 30 inputs) will be calculated, but the main analysis includes
only those inputs that all 15 networks classify correctly, reporting the number of
such inputs. Together, the records identify the computer used, the version of the
code, and the software environment for each result.

Computation is performed on 15 school laboratory computers (one network per computer), with
distribution controlled by the Raspberry Pi coordinator. Identical locked package versions are
installed on each computer, and each computer runs a self-check including the
tests and a training run to allow comparisons between computers. The coordinator will only
accept computers whose code and settings match its own. A pilot
will be run in early November using validation inputs and will check convergence,
the certificate, the attack, the frequency at which the reference resolves, and
runtime performance. Settings will be finalized on November 13, while the data collection will take
place during the week of November 16, with the week of November 30 as a reserve.

Inputs in one network are not statistically independent; thus, the network itself is
the unit of analysis. The median certificate error and the median
attack error for each network (based on interval midpoints) are calculated
over resolved instances. The five medians for each architecture are considered as
separate points for comparison. Significance tests are not planned because there are only five
networks per architecture. The report also includes clean accuracy, $\rstar$, the frequency with
which the reference resolves and the attack fails, and the corresponding runtimes.
